\chapter{Source Language}
\label{ch:source-language}

\section{Lines, Whitespace and Comments}

A source file is a sequence of instructions, directives and labels separated by whitespace. By convention each is
written on its own line, with directives and instructions indented under the label they follow. Whitespace, including
blank lines, is otherwise insignificant.

A comment starts with a semicolon and runs to the end of the line. It may follow an instruction on the same line.

% qa:assemble
\begin{lstlisting}
; a whole-line comment
    MOVI r1, #1        ; a trailing comment
\end{lstlisting}

\section{Case}

Mnemonics, condition codes, shift types and directives are written in uppercase. Register names are written in
lowercase. The status update letters inside square brackets (Section~\ref{sec:status-update-syntax}) may be either
case. Label and constant names are case-sensitive.

\section{Names}

Label names and constant names are made of letters, digits and underscores. A label is declared by writing a colon
before its name, and referred to by writing an at sign before it:

% qa:assemble
\begin{lstlisting}
:loop                 ; declare the label
    BRAN @loop        ; refer to it
\end{lstlisting}

The colon is written before the name, not after it. A label may be referred to before or after its declaration, and
from any object file linked into the same program.

\section{Numbers}

A number used as an operand value is prefixed with a hash sign. Three bases are accepted, and an underscore may be used
anywhere inside the digits to make them easier to read.

\begin{table}[H]
	\centering
	\begin{tabular}{lll}
		\toprule
		\textbf{Base} & \textbf{Written as}                    & \textbf{Examples}                               \\
		\midrule
		Decimal       & \texttt{\#} then digits, optional sign & \texttt{\#42}, \texttt{\#-3}, \texttt{\#1\_000} \\
		Hexadecimal   & \texttt{\#0x} then hex digits          & \texttt{\#0x2A}, \texttt{\#0xFF\_00}            \\
		Binary        & \texttt{\#0b} then binary digits       & \texttt{\#0b0010\_1010}                         \\
		\bottomrule
	\end{tabular}
	\caption{Number forms}
	\label{tab:number-forms}
\end{table}

A negative decimal number is stored in two's complement. The value must fit the field it is encoded into; for example
an immediate operand of a 16-bit instruction is 16 bits wide.

Two directive operands are written without the hash: the address of \texttt{.ORG} is a bare decimal or
\texttt{0x}-prefixed hexadecimal number (Section~\ref{sec:org}).

\section{Registers}

The general-purpose registers are \texttt{r1} to \texttt{r7}. The status register is \texttt{sr}. Each address register
pair is addressed as a whole by a single letter (\texttt{a}, \texttt{s}, \texttt{p} or \texttt{l}) and as two separate
halves: \texttt{ah} and \texttt{al}, \texttt{sh} and \texttt{sl}, \texttt{ph} and \texttt{pl}, \texttt{lh} and
\texttt{ll}. The \archmanual{} describes what each register is for, and which are privileged.

\section{Named Constants}

\texttt{.EQU} defines a named constant. The name starts with a dollar sign at the definition and at every use:

% qa:assemble
\begin{lstlisting}
    .EQU $ANSWER #42

    MOVI r1, $ANSWER       ; no '#' at the use
\end{lstlisting}

A constant stands in for a numeric value in an immediate operand or in a data directive. It cannot be used as the
address of \texttt{.ORG}.
